% !PRINTED 674
\weberpage{674}

%<<CONT
denominator equivalent to $\mathfrak{A}$ and the same system of constants, and this polygon $\mathfrak{A}'$ is
therefore dependent only on the constants and on the polygons
$\mathfrak{A}_1, \mathfrak{A}_2, \ldots, \mathfrak{A}_s$. The totality of these polygons is called a
\gesperrt{family of polygons with the basis} $\mathfrak{A}_1, \mathfrak{A}_2, \ldots, \mathfrak{A}_s$,
and the designation

\weq{9}{(\mathfrak{A}_1, \mathfrak{A}_2, \ldots, \mathfrak{A}_s)}

is used for it.

If the polygons $\mathfrak{A}_1, \mathfrak{A}_2, \ldots, \mathfrak{A}_s$ have the greatest common
 divisor $\mathfrak{M}$, then this divisor is contained in all the polygons
of the family, and is called the \gesperrt{divisor of the family}. But
in the family one can determine a polygon $\mathfrak{A}' = \mathfrak{M}\mathfrak{B}$ in such a way
that $\mathfrak{B}$ does not contain arbitrarily given points.

For if a point $\mathfrak{P}$ is contained $\mu$ times in $\mathfrak{M}$ and $\nu$ times in $\mathfrak{A}$,
and $\varrho$ is a function that vanishes at $\mathfrak{P}$ to the first order,
then in the developments (2):

\[
m = \mu - \nu,
\]

and the development of $\eta = \mathfrak{A}' : \mathfrak{A} = \mathfrak{M}\mathfrak{B} : \mathfrak{A}$ reads

\[
\eta = e\varrho^m + \sigma\varrho^{m+1},
\]

in which

\[
e = e_1c_1 + e_2c_2 + \cdots + e_sc_s.
\]

If, then, one chooses the $c_i$ so that $e$ does not vanish, then
$\mathfrak{P}$ is not contained in $\mathfrak{B}$.

A family that has a divisor $\mathfrak{M}$ is called an \gesperrt{improper
 family}.

If among the functions $\eta_1, \eta_2, \ldots, \eta_s$ there subsists a relation
of the form:

\weq{10}{c_1\eta_1 + c_2\eta_2 + \cdots + c_s\eta_s = 0,}

in which the constants $c_1, c_2, \ldots, c_s$ are not all $= 0$, then the
 same relation

\weq{11}{c_1\eta_1' + c_2\eta_1' + \cdots + c_s\eta_s' = 0}
% [sic] printed c_2\eta_1' as the second term, should be c_2\eta_2': (11) is (10) multiplied by \xi, with \xi\eta_2 = \eta_2' by (7) and as in (8)

subsists between the functions (7). We then call these functions
\gesperrt{linearly dependent}. Correspondingly, the polygons
of the family $\mathfrak{A}_1, \mathfrak{A}_2, \ldots, \mathfrak{A}_s$ are also called linearly dependent. One can then express one
of the functions, say $\eta_s$, linearly by the others,
and the family $(\eta_1, \eta_2, \ldots, \eta_s)$ is identical with $(\eta_1, \eta_2, \ldots, \eta_{s-1})$.
The same holds for the families of polygons $(\mathfrak{A}_1, \mathfrak{A}_2, \ldots, \mathfrak{A}_s)$ and $(\mathfrak{A}_1,
\mathfrak{A}_2, \ldots, \mathfrak{A}_{s-1})$. By repeated application of this reduction
one can represent every family by a \gesperrt{linearly independent or
irreducible basis}. Two irreducible bases of one
%CONT>>
